\subsection{A longer horizon}
\label{sec:longT}
\paragraph{Question and design} The validated action at $T=1000$ iterations is $4.89\%$ above the best known solutions in B30, so the host is not competitive in absolute terms. We repeated the benchmark in B30, the richest space, with $T=10{,}000$ iterations on the same $280$ tuning and $700$ test instances. The validated action of each class was chosen again at the long horizon, so that the static competitor has the horizon of the test; the seven learners keep their frozen hyperparameters; the seeds are independent of those of the short horizon. The six predictions $L_1$--$L_6$ and the hashes of the code and of the validated actions were recorded before the test run.

\paragraph{Results} The mean gap of the validated action to the best known solution falls to $2.69\%$ ($7.42\%\to4.31\%$ for random choice and $7.91\%\to4.23\%$ for the standard Q-Learning; $L_1$ held). The headroom shrinks with the horizon, as in the synthetic batteries: the mean gain of the validated action over random choice falls from $2.26$ to $1.50$ points ($L_5$), and it remains beyond the margin in all seven types ($L_3$; Table~\ref{tab:long}). The validated action again beats all $49$ learner--type pairs, and no learner beats it beyond the margin ($L_2$). The best learner of each type recovers from $12\%$ to $39\%$ of the validated gain ($L_6$ held in all seven types). The online learners fare better at the long horizon than at the short one in one respect only: the ALNS roulette gains beyond the margin in six of the seven types (up to $1.18$ points on $U[10,100]$, a third of the validated gain, as it did at $T=1000$), the episodic published design in four and the FRRMAB-style design in five. The standard Q-Learning gains in two types ($0.07$ and $0.45$ points) against a prediction of at most one ($L_4$ failed); five of the six predictions held. The order-invariant Q-Learning (Entry~263, registered before it was run) shows that these two gains are the same artefact: it gains beyond the margin in none of the seven types, is worse than random choice in six (by $0.09$ to $0.90$ points) and level with it in the seventh, and the validated action beats it in all seven types. It evaluates $1.10$ times as many solutions as random choice, against $0.72$ for the standard learner. Four of the five predictions of that test held; the miss was the evaluation ratio, registered as $0.9$--$1.1$ and observed just above it.

% Auto-generated by analysis/make_table_journal.py
\begin{table}[t]
\caption{Standard benchmark, B30, longer horizon: gain over random choice (points of makespan, positive = better; $100$ test instances per type, three seeds) of the validated action at $T=1000$ and at $T=10{,}000$ (re-chosen at the long horizon), of the standard Q-Learning (ql), the ALNS roulette (alns) and of the best of the seven learners at $T=10{,}000$ (best; the learner is named), and the gap of the validated action to the best known solution (\%). Bold: gain over random choice confirmed (Holm within learner, seven types) and not equivalent to zero at the margin $\pm0.05$.}
\label{tab:long}
\centering
{\scriptsize\setlength{\tabcolsep}{4pt}
\resizebox{\textwidth}{!}{%
\begin{tabular}{@{}lrrrrrlrr@{}}
\toprule
 & \multicolumn{2}{c}{validated gain} & \multicolumn{4}{c}{learners at $T=10{,}000$} & \multicolumn{2}{c}{validated gap (\%)} \\
\cmidrule(lr){2-3}\cmidrule(lr){4-7}\cmidrule(lr){8-9}
Type & $T=10^3$ & $T=10^4$ & ql & alns & \multicolumn{2}{c}{best} & $T=10^3$ & $T=10^4$ \\
\midrule
$U[1,100]$ & 5.08 & 2.69 & 0.12 & \textbf{0.95} & \textbf{0.95} & alns & 6.8 & 3.8 \\
$U[10,100]$ & 4.15 & 3.52 & \textbf{0.45} & \textbf{1.18} & \textbf{1.18} & alns & 9.9 & 5.5 \\
$U[100,120]$ & 0.51 & 0.26 & $-$0.04 & $-$0.04 & 0.04 & ucb & 2.1 & 1.5 \\
$U[100,200]$ & 2.93 & 2.39 & 0.09 & \textbf{0.16} & \textbf{0.94} & iq & 4.8 & 1.9 \\
$U[1000,1100]$ & 0.40 & 0.21 & \textbf{0.07} & \textbf{0.03} & \textbf{0.08} & iq & 0.5 & 0.3 \\
correlated jobs & 0.50 & 0.54 & $-$0.18 & \textbf{0.07} & \textbf{0.07} & alns & 3.3 & 2.3 \\
correlated machines & 2.22 & 0.90 & $-$0.13 & \textbf{0.26} & \textbf{0.26} & alns & 6.9 & 3.7 \\
\midrule
mean & 2.26 & 1.50 & 0.06 & 0.37 & 0.50 & & 4.9 & 2.7 \\
\bottomrule
\end{tabular}}}
\end{table}


\subsection{A policy trained offline with a neural network}
\label{sec:neural}
\paragraph{Question and design} The offline selector of Section~\ref{sec:selector} is a table indexed by the class; the closest relative in the algorithm-configuration literature is a policy that maps the state of the search to the configuration and is trained over many instances before it is used~\cite{BiedenkappEtal2020,EimerEtal2021}. We built one for the IG of unrelated machines: a network with one hidden layer ($16$ tanh units) maps eight scale-free features (the fraction of the horizon spent, the iterations without improvement, the distance to the lower bound, the load imbalance, $\ln n$, $\ln m$, a moving average of recent improvements and a bias) to a softmax over the actions. The weights are trained by evolution strategies~\cite{SalimansEtal2017} (antithetic perturbations, rank-based fitness shaping, Adam) to minimise the ratio of the makespan to the best known value over the $280$ tuning instances, and frozen at test time: nothing is learned during a run. A control, \textsc{nn-const}, has no features, so it is a constant mixture of the actions, trained with the same budget; it separates what the state adds from what a good mixture adds. Training took $138{,}240$ IG runs in B30 and $115{,}200$ in B15 for the full policy ($69{,}120$ for each control), against $16{,}800$ and $8{,}400$ runs for the validated action; this cost is not charged to the policy in the comparison. The standard benchmark, B15 and B30, $700$ test instances and $T=1000$; the predictions $N_1$--$N_5$ and the hashes were recorded before the test run.

\paragraph{Results} The policy gains over random choice beyond the margin in all $14$ cells ($N_1$), more than the best of the seven online learners in all $14$ ($N_3$), and it never beats the validated action beyond the margin ($N_2$), which beats it, confirmed, in $11$ of the $14$ cells. Averaged over the $14$ cells the validated action gains $1.95$ points, the policy $1.31$ ($67\%$; $N_4$), the constant mixture $1.17$ ($60\%$) and the best online learner of each cell $0.35$ ($18\%$; Table~\ref{tab:neural}). The state helps, but modestly: the policy beats the constant mixture beyond the margin in $6$ of the $14$ cells ($N_5$ held, narrowly; all six are in B30, by $0.04$ to $0.50$ points) and in none of the B15 cells. All five predictions held. So the knowledge acquired before the run, even without a class, is worth about $3.7$ times what the best online learner finds during it; most of that value is a good mixture of actions, not a decision that depends on the state; and a static action chosen per class (a class the policy does not see) on about a tenth of the training effort is still better.

% Auto-generated by analysis/make_table_journal.py
\begin{table}[t]
\caption{A neural policy trained offline on the standard benchmark: gain over random choice (points of makespan, positive = better; $100$ test instances per cell, $T=1000$, three seeds) of the validated action (val), of the policy with state and instance features (nn), of a constant mixture of actions trained with the same budget (nn\_const) and of the best of the seven online learners in each cell (online; the learner is named). Bold: nn better than nn\_const, confirmed (Holm over the $14$ cells) and not equivalent to zero at the margin $\pm0.05$. nn never beats val beyond the margin.}
\label{tab:neural}
\centering
{\scriptsize\setlength{\tabcolsep}{4pt}
\resizebox{\textwidth}{!}{%
\begin{tabular}{@{}lrrrrlrrrrl@{}}
\toprule
 & \multicolumn{5}{c}{B15} & \multicolumn{5}{c}{B30} \\
\cmidrule(lr){2-6}\cmidrule(lr){7-11}
Type & val & nn & nn\_const & \multicolumn{2}{c}{online} & val & nn & nn\_const & \multicolumn{2}{c}{online} \\
\midrule
$U[1,100]$ & 4.22 & 2.03 & 2.00 & 0.54 & fr & 5.08 & \textbf{3.16} & 2.77 & 0.78 & qo \\
$U[10,100]$ & 3.27 & 2.23 & 2.07 & 0.66 & iq & 4.15 & \textbf{3.12} & 2.63 & 0.61 & qo \\
$U[100,120]$ & 0.60 & 0.53 & 0.55 & 0.25 & iq & 0.51 & \textbf{0.42} & 0.33 & 0.10 & qo \\
$U[100,200]$ & 1.56 & 1.23 & 1.19 & 0.32 & iq & 2.93 & \textbf{1.99} & 1.63 & 0.39 & alns \\
$U[1000,1100]$ & 0.43 & 0.38 & 0.39 & 0.14 & fr & 0.40 & \textbf{0.30} & 0.26 & 0.09 & qo \\
correlated jobs & 0.44 & 0.27 & 0.26 & 0.02 & fr & 0.50 & 0.16 & 0.08 & 0.13 & alns \\
correlated machines & 0.94 & 0.83 & 0.79 & 0.26 & fr & 2.22 & \textbf{1.68} & 1.41 & 0.63 & alns \\
\midrule
mean, 14 cells & \multicolumn{10}{l}{val 1.95; nn 1.31; nn\_const 1.17; best online 0.35} \\
\bottomrule
\end{tabular}}}
\end{table}

